% =============================================================
%  Relatório — Trabalho M1: Classificação Multiclasse com MLP
%  Disciplina: Inteligência Artificial
%  Alunos: Gabriel Schaldach Morgado, Eduardo Lechinski Ramos
% =============================================================
\documentclass[12pt,a4paper]{article}

% ---------- Pacotes ----------
\usepackage[utf8]{inputenc}
\usepackage[T1]{fontenc}
\usepackage[brazilian]{babel}
\usepackage{geometry}
\usepackage{graphicx}
\usepackage{booktabs}
\usepackage{longtable}
\usepackage{hyperref}
\usepackage{amsmath}
\usepackage{array}
\usepackage{xcolor}
\usepackage{caption}
\usepackage{subcaption}
\usepackage{float}
\usepackage{lmodern}
\usepackage{parskip}

\geometry{margin=2.5cm}

\hypersetup{
    colorlinks=true,
    linkcolor=blue,
    urlcolor=blue,
    citecolor=blue
}

% ---------- Cabeçalho ----------
\title{\textbf{Trabalho Prático M1}\\
       \large Classificação Multiclasse com Redes Neurais\\
       \normalsize Dry Bean Dataset}
\author{Gabriel Schaldach Morgado \and Eduardo Lechinski Ramos}
\date{\today}

% =============================================================
\begin{document}
\maketitle
\tableofcontents
\newpage

% =============================================================
\section{Introdução}
% =============================================================

Este relatório descreve o desenvolvimento de um classificador multiclasse
baseado em Multilayer Perceptron (MLP) para identificar sete variedades de
feijão a partir de características morfológicas extraídas de imagens dos grãos.
O dataset utilizado é o \textit{Dry Bean Dataset}
(\url{https://archive.ics.uci.edu/dataset/602/dry+bean+dataset}), criado por
Koklu e Ozkan (2020).

% =============================================================
\section{Análise Exploratória de Dados}
% =============================================================

\subsection{Inspeção Inicial}

\textbf{Q1 — Dimensões do dataset:}
O dataset possui 13.611 linhas e 17 colunas (16 features numéricas + 1 coluna
alvo \texttt{Class}), em conformidade com a documentação do UCI.

\textbf{Q2 — Valores ausentes e tipos de features:}
Não há valores ausentes.  Das 16 features, 2 são inteiras (\texttt{Area} e
\texttt{ConvexArea}) e 14 são contínuas; a variável alvo é categórica.

\textbf{Q3 — Desbalanceamento de classes:}
A classe BOMBAY possui apenas 522 amostras, enquanto DERMASON possui 3.546.
Esse desequilíbrio pode levar a rede a favorecer classes majoritárias,
gerando alta acurácia global mas baixo recall para classes raras.
A acurácia sozinha pode mascarar esse problema: um modelo que sempre prediz
DERMASON já obteria cerca de 26\% de acurácia "de graça".

\subsection{Duplicatas}

\textbf{Q4 — Registros duplicados:}
Foram encontrados 68 registros duplicados no dataset completo.
O impacto é mínimo dado o volume total (< 0,5\%), mas foram removidos
para evitar vazamento de informação entre treino e teste.

\subsection{Distribuição das Features}

\textbf{Q5 — Features assimétricas com cauda longa:}
\texttt{Area}, \texttt{Perimeter}, \texttt{MajorAxisLength},
\texttt{MinorAxisLength}, \texttt{Eccentricity}, \texttt{ConvexArea},
\texttt{EquivDiameter} e \texttt{Roundness} apresentam distribuições
assimétricas evidentes nos histogramas.

\textbf{Q6 — Efeito da transformação logarítmica:}
A transformação \texttt{log1p} aproximou essas distribuições da normalidade,
centralizando a média e atenuando os valores extremos.
Isso beneficia o treinamento ao reduzir a influência de outliers nos gradientes
e garantir que o StandardScaler opere sobre uma distribuição mais simétrica.

\subsection{Correlação de Pearson}

\textbf{Q7 — Pares com correlação acima de 0,90:}
\begin{itemize}
    \item \texttt{Area} e \texttt{Perimeter} ($r > 0{,}95$)
    \item \texttt{ShapeFactor3} e \texttt{Compactness} ($r > 0{,}92$)
    \item \texttt{MinorAxisLength} e \texttt{ConvexArea} ($r > 0{,}93$)
\end{itemize}

\textbf{Q8 — Multicolinearidade:}
Multicolinearidade ocorre quando variáveis independentes possuem relações
lineares exatas ou aproximadas entre si.  Features altamente correlacionadas
carregam informação redundante.  Contudo, \textbf{não é obrigatório} removê-las:
se a feature é relevante para a tarefa, sua remoção pode introduzir erro de
especificação.

\subsection{Boxplots por Classe}

\textbf{Q9 — Features selecionadas para o modelo:}
\textit{Area}, \textit{MajorAxisLength}, \textit{MinorAxisLength},
\textit{AspectRation}, \textit{Eccentricity}, \textit{Extent},
\textit{Solidity}, \textit{Roundness}, \textit{Compactness},
\textit{ShapeFactor1}, \textit{ShapeFactor2} e \textit{ShapeFactor4}.
As features \texttt{ShapeFactor3}, \texttt{ConvexArea}, \texttt{Perimeter}
e \texttt{EquivDiameter} foram excluídas por alta correlação com outras.

\textbf{Q10 — Features que melhor separam as classes visualmente:}
\texttt{ShapeFactor1} e \texttt{Area} apresentaram a maior separação visual
nos boxplots, com distribuições bem distintas por classe.

\textbf{Q11 — Classes com maior sobreposição:}
BARBUNYA e HOROZ apresentaram sobreposição considerável em várias features.
Espera-se que o modelo confunda essas classes com mais frequência.

% =============================================================
\section{Pré-processamento}
% =============================================================


\subsection{Normalização das Features}

\textbf{Q12 — Por que não ajustar o scaler nos dados de teste?}
O scaler foi ajustado exclusivamente nos dados de treino.  Usar dados de teste
no \texttt{fit} introduziria \textit{data leakage}: os parâmetros de
escalonamento incorporariam informações do conjunto de avaliação.

\textbf{Q13 — Scaler diferente em produção:}
Se um scaler diferente for aplicado na inferência, as entradas chegam à rede
em escala distinta da vista no treinamento.  Os pesos da rede foram otimizados
para uma distribuição específica de entradas, fazendo o modelo ter predições incorretas.

\subsection{Pesos de Classe}

\textbf{Q14 — BOMBAY recebeu o maior peso:}
Sim.  Com apenas 522 amostras, BOMBAY obteve o maior peso (\(\approx 3{,}7\)),
inversamente proporcional à sua frequência no dataset.

\textbf{Q15 — Impacto dos pesos na função de perda:}
Um erro em BOMBAY gera uma penalização aproximadamente 4× maior do que um
erro em DERMASON.  Isso força a rede a dar mais atenção a exemplos raros
durante a descida de gradiente.

\subsection{Divisão Treino / Validação / Teste}

\textbf{Q16 — Importância do \texttt{stratify}:}
O \texttt{stratify=y} garante que cada partição mantenha a proporção original
de cada classe.  Sem ele, um split puramente aleatório poderia concentrar
quase todos os exemplos de BOMBAY (522 no total) em um único.

\textbf{Q17 — Dados de teste apenas na avaliação final:}
Usar dados de teste para escolher arquitetura ou hiperparâmetros junta eles
implicitamente ao processo de desenvolvimento.  A métrica obtida deixa de
ser uma estimativa imparcial do desempenho em dados realmente novos.

% =============================================================
\section{Construção, Treinamento e Experimentos}
% =============================================================

\subsection{Arquitetura e Decisões de Design}

\textbf{Q18 — Softmax com 7 neurônios na saída:}
A função softmax transforma as ativações finais em uma distribuição de
probabilidade que soma 1,0, com um neurônio por classe.

\textbf{Q19 — \texttt{sparse\_categorical\_crossentropy}:}
Essa função de perda opera sobre rótulos inteiros e não binários diferente de \texttt{binary\_crossentropy} que é restrita à classificação
binária. \texttt{categorical\_crossentropy} exigiria rótulos em outros formatos.

\textbf{Q20 — Dropout na camada de saída:}
Não.  O Dropout é aplicado apenas nas camadas ocultas.  Aplicar ele na camada
de saída desativaria neurônios responsáveis por classes inteiras.

\subsection{Tabela Comparativa de Experimentos}

% PREENCHA OS VALORES COM OS RESULTADOS REAIS DO SEU NOTEBOOK
\begin{table}[H]
\centering
\caption{Resultados dos experimentos no conjunto de teste}
\label{tab:experimentos}
\resizebox{\textwidth}{!}{%
\begin{tabular}{@{}llllrrrl@{}}
\toprule
\textbf{Exp.} & \textbf{Descrição} & \textbf{Features} &
\textbf{Acurácia} & \textbf{F1-macro} &
\textbf{Observações} \\
\midrule
A  & Sem Dropout (baseline)              & Todas (16)         & .9199 & .9305  & Overfitting visível nas curvas \\
A2 & Com Dropout (0.3 / 0.2)             & Todas (16)         & .9225 & .9344  & Gap treino-val reduzido (não chegou a reproduzir o overfitting) \\
B  & A2 + class\_weight                  & Todas (16)         & .9240 & .9355  & Não demonstrou diferença e aumentou o gap de val-loss \\
C  & Rede maior [256-128-64]             & Todas (16)         & .9221 & .9347  & redução de overfit e rápida convergência \\
D  & Inicializador RandomNormal          & Todas (16)         & .9254 & .9365  & Convergência mais lenta \\
E  & Features selecionadas na AED        & Selecionadas (12)  & .9232 & .9353  & Maior confusão entre classes similares \\
F  & Melhor configuração                 & Todas (16)         & .9266 & .9380  & Melhor F1-macro geral \\
\bottomrule
\end{tabular}}
\end{table}


% =============================================================
\section{Avaliação e Análise dos Resultados}
% =============================================================

\subsection{Matriz de Confusão — Melhor Modelo (Experimento F)}

% Insira a figura gerada pelo notebook:
% \begin{figure}[H]
%   \centering
%   \includegraphics[width=0.75\textwidth]{confusion_matrix_F.png}
%   \caption{Matriz de confusão do Experimento F no conjunto de teste.}
%   \label{fig:cm}
% \end{figure}

\subsection{Análise Qualitativa}

\textbf{Q23 — Classe com menor F1-score:}
Esperamos que SIRA ou BARBUNYA apresentem os menores F1-scores, o que reflete a sobreposição vista nos boxplots.


\textbf{Q24 — Classes confundidas com mais frequência:}
A matriz de confusão mostra mais entre SIRA e DERMASON,
o que tinha sido identificado na
sobreposição na AED.

\textbf{Q25 — BOMBAY e o efeito de class\_weight:}
A classe BOMBAY foi bem identificada em todos os resultados. O class\_weight nesse caso não teve nenhuma mudança expressiva nessa classe.

\textbf{Q28 — Desbalanceamento e class\_weight:}
Não houve muitos sinais de efeito no desbalanceamento com os parametros usados, não houve muita mudança nos resultados. As classes desbalanceadas possuem propriedades bem definidas.

\textbf{Q29 — Melhor experimento:}
O Experimento F produziu o melhor F1-macro porem foram feitos por muitas iterações de teste. O melhor expiremento antes dos testes e iteração foi o experimento C

\textbf{Q30 — Dropout (Exp A vs.\ Exp A2):}
O dropout reduziu muito o overfit, principalmente quando olhamos para o experimento sem early stopping. No experimento A2 a curva de val\_loss se mostrou muito mais comportada.

\textbf{Q31 — Features selecionadas (Exp E):}
A remoção de features redundantes não produziu degradação significativa. Na verdade aumentou confusão entre algumas classes similares.

\textbf{Q32 — Justificativa do Experimento F:}
A arquitetura do experimento F foi baseada em testes. Rodando cada modelo retirando e colocando parametros e vendo o que melhora em cada um, a maior camada e o dropout reduziram consideravelmente o overfit da rede. Diminuimos a taxa de aprendizado por conta dos saltos que podemos ver no decaimento da acurácia dos modelos E por fim 

% =============================================================
\section{Conclusão}
% =============================================================

O Experimento F demonstrou que a combinação de arquitetura mais profunda,
Dropout e uma taxa de aprendizado menor produz o melhor acurácia global e F1-macro. Porém os experimentos foram consistentes retirando os reultados de overfit como o experimento A e A2.
A análise exploratória foi importante para identificar as classes e peculiaridades que deveriamos ter tomado cuidado ao produzir o modelo. Entender o dataset foi essencial para a melhoria dos parametros.

O uso de \texttt{class\_weight} mostrou-se pouco eficaz, porém reduziu mesmo que minimamente o overfit e manteve o modelo mais estável durante os testes.

% =============================================================
\end{document}
